% Ad Familiares 6.1 --- headnote
% ad-familiares-06-01, late intercalary month 46 BC, Rome

\textit{Cicero to A. Manlius Torquatus, written at Rome at the
end of the second intercalary month of 46 BC (Perseus:
\textit{Romae ex.~m.~intercal.~post.~a.~708 (46)}). Torquatus, a
Pompeian who had not been pardoned with the rest after Pharsalus,
was still living away from Rome in continued political exile.
The letter is an attempt at consolation written by one defeated
former Pompeian to another, with the additional weight that
Cicero --- restored to the city but not to public life --- was
trying to argue his friend back from despair toward the form of
Stoic composure he was at the same time arguing himself toward.}

\textit{The letter is philosophically pitched in a way that
anticipates the great consolatory works of the following year:
the lost \textit{Consolatio} for Tullia, the \textit{Tusculan
Disputations}, parts of \textit{De Officiis}. The central
argument of \S\S 3--4 --- that what we owed the commonwealth was
duty, not the demand for a particular outcome; that, the outcome
having gone against us, we have nothing to be ashamed of, only
something to bear --- is the doctrine on which Cicero will lean
for the rest of his philosophical period. The reasoning that
the defeated man must not feign surprise at what was always among
the possible outcomes (``We were not so deranged as to suppose
victory was assured to us'') is offered in the voice of a man
who has been over the same ground himself.}

\textit{In \S 6 Cicero notes that Torquatus has Servius Sulpicius
Rufus near him --- the same Servius who in mid-March of the
following year will write to Cicero on the death of Tullia
(\textit{Fam.}~4.5). The valedictory section returns to the
private register: those to whom Cicero owed most ``the chances of
this war have torn from me,'' and what remains of his counsel and
his effort is owed to Torquatus and his children. The closing
balances --- ``under some form of commonwealth you will be the
man you ought to be, or with the commonwealth ruined your case
will be no worse than every other man's'' --- carry the
characteristic Ciceronian antithesis of the period: there are
two futures, and the philosophical man is ready for either.}
